\documentclass[10pt, a4paper]{article}
\usepackage[a4paper,
  total={6in, 8in},
  left=0.55in,
  right=0.55in,
  top=0.70in,
  bottom=0.25in,
  nohead
]{geometry}

% Essential packages for academic formatting
\usepackage[utf8]{inputenc}
\usepackage[T1]{fontenc}
\usepackage[margin=1in]{geometry} % Standard margins
\usepackage{titling}
\usepackage{enumitem} % For better list control
\usepackage{parskip}  % Adds space between paragraphs for readability
\usepackage{textcomp} % For the Euro symbol

% Title Setup
\title{\textbf{Research Statement}}
\author{}
\date{}

\begin{document}

\maketitle
\vspace{-2em} % Adjust vertical space to bring text closer to title

\section*{Research Vision: Grounded Causal Intelligence for Net-Zero Manufacturing \& Infrastructure}

My research establishes a framework for \textbf{Grounded Causal Intelligence (GCI)}---a symbiotic loop where AI agents autonomously design physical experiments to correct physics-based models. While current ``Digital Twins'' often function as static observers, my work transforms them into adaptive agents capable of ``Zero-Shot Planning.'' 

This transition from reactive data logging to proactive causal reasoning is critical for two strategic priorities at the University of Birmingham: \textbf{Net-Zero Manufacturing} and \textbf{Resilient Transportation Infrastructure}.

\section*{Research Pillars \& Strategic Alignment}

\subsection*{1. Adaptive Digital Twins for Zero-Defect Manufacturing}
The primary application of GCI is to eliminate the ``physical residual''---the stochastic gap between simulation and reality. By integrating multi-modal sensor data (acoustic, thermal, visual) with Causal AI, my models autonomously identify the root cause of defects during production.

\begin{description}
    \item[Application:] I will utilize the Department’s advanced manufacturing facilities to scale my current work on Laser Powder Bed Fusion (LPBF) and Wire Arc Additive Manufacturing (WAAM).
    \item[Impact:] This supports the School’s \textit{Sustainable Manufacturing} theme by reducing material waste and energy consumption associated with trial-and-error qualification.
\end{description}



\subsection*{2. AI-Driven Resilience for Transportation Systems (UKRRIN Alignment)}
The mathematical core of my work—using sensor data to update physics models in real-time—is directly transferable from manufacturing to infrastructure maintenance.

\begin{description}
    \item[Proposal:] I aim to collaborate with the \textbf{UK Rail Research and Innovation Network (UKRRIN)} to apply GCI to predictive maintenance for rolling stock and rail infrastructure.
    \item[Methodology:] By treating a rail component similarly to a manufacturing workpiece, ``physical residuals'' in vibration or acoustic data can be causally linked to early-stage fatigue or structural degradation, enabling ``preventive intervention'' rather than ``reactive repair.''
\end{description}

\subsection*{3. Data-Centric Engineering for Buried Infrastructure}
Aligning with the \textbf{National Buried Infrastructure Facility (NBIF)}, I propose extending my ``Context Extractor'' algorithms (currently used for 3D CAD analysis) to interpret subterranean sensor data. This would enable the automated generation of ``underground digital twins'' that update dynamically as soil conditions or structural loads change.

\section*{Funding Strategy}
My objective is to secure sustainable funding to support a research group of 3--4 PhD students and postdoctoral researchers.

\begin{itemize}[leftmargin=*]
    \item \textbf{Primary Target: EPSRC New Investigator Award.} I will submit a proposal to port the ``Grounded Causal Intelligence'' engine to UK high-value manufacturing sectors.
    \item \textbf{Collaborative Grants:} Leveraging my role in the AI4AM consortium (Horizon Europe, under review) and AutoCAM (FWO), I will target Horizon Europe Cluster 4 and Innovate UK calls focused on ``Digital \& Emerging Technologies.''
\end{itemize}

\section*{Expected Outcomes}

\begin{itemize}[leftmargin=*]
    \item \textbf{REF 2029:} I commit to publishing 2--3 high-impact papers annually in journals such as \textit{Nature Communications}, \textit{Additive Manufacturing}, and \textit{IEEE Transactions on Industrial Informatics}, ensuring a strong contribution to the School’s REF profile.
    \item \textbf{Industrial Impact:} I will continue my practice of releasing open-source tools (e.g., SmartECM) to facilitate rapid industrial uptake of academic research.
\end{itemize}

\end{document}